% !TeX root = ../main.tex
\section{Программная реализация моделирования}
Расчет аналитических зависимостей и дискретное моделирование системы ЧАП выполнены на языке Python с использованием библиотек \texttt{numpy} и \texttt{matplotlib}.

\subsection{Функции аналитического расчета установившегося режима}
\begin{lstlisting}[language=Python]
import numpy as np
import matplotlib.pyplot as plt

alpha = 1.0
T = 0.01
f0 = 1602e6
omega0 = 2.0 * np.pi * f0
c = 3e8
k_dop = omega0 / c

def calc_N0_tilde(q_lin):
    return (2.0 / (q_lin * T**2)) * (1.0 + 1.0 / (2.0 * q_lin * T))

def calc_S_xi(sigma_a):
    return 2.0 * (sigma_a**2) * alpha * (k_dop**2)

def steady_state_params(q_dBHz, sigma_a):
    q_lin = 10.0**(0.1 * q_dBHz)
    N0_t = calc_N0_tilde(q_lin)
    S_xi = calc_S_xi(sigma_a)

    D11 = (alpha * N0_t / 2.0) * (
        np.sqrt(1.0 + 2.0 * np.sqrt(S_xi / (alpha**2 * N0_t))) - 1.0
    )
    sigma_omega = np.sqrt(D11)

    K1 = 2.0 * D11 / N0_t
    K2 = (K1**2) / 2.0

    num = (K1**2) * (alpha + K1) + (alpha * K1 + K2)**2
    den = 4.0 * (alpha + K1) * (alpha * K1 + K2)
    dF_fll = num / den

    return sigma_omega, dF_fll, D11
\end{lstlisting}

\subsection{Цикл дискретного фильтра Калмана}
\begin{lstlisting}[language=Python]
F = np.array([[1.0, T], [0.0, 1.0 - alpha * T]])
G = np.array([[0.0], [alpha * T]])
H = np.array([[1.0, 0.0]])
Q_d = (G @ G.T) * (sigma_xi_d**2)
R_d = sigma_n_d**2

x_true = np.array([100.0, 100.0])
x_est = np.array([0.0, 0.0])
D = np.array([[34.0**2, 0.0], [0.0, 340.0**2]])

for k in range(N_steps):
    omega_true_hist[k] = x_true[0]
    omega_est_hist[k] = x_est[0]
    D11_hist[k] = D[0, 0]

    x_true = F @ x_true + np.array([0.0, alpha * T * xi_noise[k]])
    y_meas = x_true[0] + n_noise[k]

    x_pred = F @ x_est
    D_pred = F @ D @ F.T + Q_d

    S_inn = D_pred[0, 0] + R_d
    K = D_pred[:, 0] / S_inn

    innov = y_meas - x_pred[0]
    x_est = x_pred + K * innov
    D = (np.eye(2) - np.outer(K, H)) @ D_pred
\end{lstlisting}
